\subsection{Variadic Positional Parameters}

A variadic positional parameter (\texttt{*args}) collects remaining positional arguments into a heap-allocated \texttt{tuple}. The star in the definition packs; the star at the call site unpacks an iterable into separate positional arguments.

\subsubsection{Variable-Length Argument Packing: The Mechanics of \texttt{*args} and \texttt{**kwargs}}

When a signature includes \texttt{*args}, CPython must allocate a fresh \texttt{tuple} at the call site to hold overflow positional references. When it includes \texttt{**kwargs}, it allocates a fresh \texttt{dict} for unbound keyword pairs. Both are ordinary heap objects with full allocation cost.

Contrast this with the Vectorcall protocol (Python 3.8+): for many C-accelerated and simple Python calls, arguments pass through a contiguous C array of \texttt{PyObject*} pointers with a method flag (\texttt{METH\_VECTORCALL}). No intermediate tuple or dict is constructed. Signatures that force \texttt{*args}/\texttt{**kwargs}---or wrappers that use them---disable this fast path and pay explicit container allocation on every invocation.

\begin{codebox}[]{python}{Variable-Length Argument Packing: The Mechanics of *args and **kwargs}{var01}
def flexible(*args, **kwargs):
    return len(args), len(kwargs)

def fixed(a, b):
    return a + b

print(flexible(1, 2, 3, x=4))
print(fixed(1, 2))
\end{codebox}


\begin{iobox}{Expected Output}
\texttt{(3, 1)} \\
\texttt{3} \\
\end{iobox}




Decorators and forwarding wrappers commonly use \texttt{def wrapper(*args, **kwargs): return fn(*args, **kwargs)}. This preserves call shapes but forces tuple/dict packing at the wrapper boundary even when the underlying callable could have accepted a vectorcall layout.

\subsubsection{Argument Unpacking Operations: Expanding Sequences Across Function Call Boundaries with \texttt{*}}

Call-site \texttt{*iterable} expands an iterable into positional arguments before binding. The iterable must be iterable; element count must match the receiving signature unless a variadic parameter absorbs overflow. Unpacking passes object references, not copies.

\begin{verbatim}
def show_line(name, score, mood):
    print(f"{name}, {score}, {mood}")

data = ["Homer", 42, "hungry"]
show_line(*data)
\end{verbatim}

Possible output:

\begin{verbatim}
Homer, 42, hungry
\end{verbatim}

Definition-side \texttt{*args} and call-side \texttt{*} are mirror operations: one packs many references into one tuple; the other expands one iterable into many arguments.
